% GENERATED FILE. Edit canonical lessons, then run npm run generate:books.
% canonical-source-hash: fnv1a64:407346f83be16326
% canonical-lessons: UR-C56-mirch, UR-C56-makkhan, UR-C56-panir, UR-C56-shahd, UR-C56-angur

\chapter{Chilli, Butter, Cheese, Honey, Grapes}
\label{ch:ur-chilli-butter-cheese-honey-grapes}
\hlchaptermodality{56}

\begin{chapteropening}
\textbf{By the end of this chapter:} \emph{I can say chilli, pepper, butter, cheese, honey and grapes.}

Complete the last lesson of chapter 56: I can say chilli, pepper, butter, cheese, honey and grapes.
\end{chapteropening}

\section[\texorpdfstring{\ur{مرچ} (mirch)}{مرچ (mirch)}]{\ur{مرچ} (mirch) — chilli, pepper}

\label{lesson:UR-C56-mirch}



\begin{quote}
Before the new one: say the Urdu for sugar, then the Urdu for salt.
\end{quote}

\subsection*{You'll want to know: \ur{مرچ}}
\textbf{\ur{مرچ}} — \emph{mirch} — \textquotedblleft{}chilli, pepper\textquotedblright{}.

\begin{cousinweb}
\textbf{\ur{مرچ}} (\emph{mirch}) is a tadbhava, a word inherited and reshaped by sound change. It comes from \emph{*maricca}, \textquotedblleft{}peppercorn\textquotedblright{}, a form CDIAL reconstructs from its descendants (CDIAL 9875). It is a variant CDIAL reconstructs beside Sanskrit \emph{marīca}, \textquotedblleft{}peppercorn\textquotedblright{}, which it takes from an Austroasiatic language. When the chilli came to India, it took the name of black pepper.
\end{cousinweb}

\begin{grammarlens}[title={What you've built}]
One more word for this chapter.
\end{grammarlens}

\subsection*{Your turn}
\begin{itemize}
\raggedright
  \item \emph{Say it:} \emph{mirch}
  \item \emph{Say it:} \emph{mirch}, once more
  \item \emph{Say it:} say \emph{mirch}
  \item \emph{Recall:} say the Urdu for sugar, then the Urdu for salt, then say \emph{mirch} again
\end{itemize}

\begin{tcolorbox}[breakable,colback=teal!4,colframe=teal!35!black,title={Before you move on}]
What does \ur{مرچ} mean? (\textquotedblleft{}Chilli, pepper\textquotedblright{}.) Say it once more.
\end{tcolorbox}
\section[\texorpdfstring{\ur{مکھن} (makkhan)}{مکھن (makkhan)}]{\ur{مکھن} (makkhan) — butter}

\label{lesson:UR-C56-makkhan}



\begin{quote}
Before the new one: say the Urdu for salt, then the Urdu for chilli.
\end{quote}

\subsection*{You'll want to know: \ur{مکھن}}
\textbf{\ur{مکھن}} — \emph{makkhan} — \textquotedblleft{}butter\textquotedblright{}.

\begin{cousinweb}
\textbf{\ur{مکھن}} (\emph{makkhan}) is a tadbhava, a word inherited from Sanskrit and reshaped by sound change. It comes from \emph{mrakṣaṇa}, \textquotedblleft{}rubbing in; ointment, oil\textquotedblright{} (CDIAL 10378). Butter was first \textquotedblleft{}the thing rubbed in\textquotedblright{}.
\end{cousinweb}

\begin{grammarlens}[title={What you've built}]
One more word for this chapter.
\end{grammarlens}

\subsection*{Your turn}
\begin{itemize}
\raggedright
  \item \emph{Say it:} \emph{makkhan}
  \item \emph{Say it:} \emph{makkhan}, once more
  \item \emph{Say it:} say \emph{makkhan}
  \item \emph{Recall:} say the Urdu for salt, then the Urdu for chilli, then say \emph{makkhan} again
\end{itemize}

\begin{tcolorbox}[breakable,colback=teal!4,colframe=teal!35!black,title={Before you move on}]
What does \ur{مکھن} mean? (\textquotedblleft{}Butter\textquotedblright{}.) Say it once more.
\end{tcolorbox}
\section[\texorpdfstring{\ur{پنیر} (panīr)}{پنیر (panīr)}]{\ur{پنیر} (panīr) — cheese}

\label{lesson:UR-C56-panir}



\begin{quote}
Before the new one: say the Urdu for chilli, then the Urdu for butter.
\end{quote}

\subsection*{You'll want to know: \ur{پنیر}}
\textbf{\ur{پنیر}} — \emph{panīr} — \textquotedblleft{}cheese\textquotedblright{}.

\begin{cousinweb}
\textbf{\ur{پنیر}} (\emph{panīr}) is a loan, usually traced to Persian \emph{panīr}, \textquotedblleft{}cheese\textquotedblright{}.
\end{cousinweb}

\begin{grammarlens}[title={What you've built}]
One more word for this chapter.
\end{grammarlens}

\subsection*{Your turn}
\begin{itemize}
\raggedright
  \item \emph{Say it:} \emph{panīr}
  \item \emph{Say it:} \emph{panīr}, once more
  \item \emph{Say it:} say \emph{panīr}
  \item \emph{Recall:} say the Urdu for chilli, then the Urdu for butter, then say \emph{panīr} again
\end{itemize}

\begin{tcolorbox}[breakable,colback=teal!4,colframe=teal!35!black,title={Before you move on}]
What does \ur{پنیر} mean? (\textquotedblleft{}Cheese\textquotedblright{}.) Say it once more.
\end{tcolorbox}
\section[\texorpdfstring{\ur{شہد} (shahd)}{شہد (shahd)}]{\ur{شہد} (shahd) — honey}

\label{lesson:UR-C56-shahd}



\begin{quote}
Before the new one: say the Urdu for butter, then the Urdu for cheese.
\end{quote}

\subsection*{You'll want to know: \ur{شہد}}
\textbf{\ur{شہد}} — \emph{shahd} — \textquotedblleft{}honey\textquotedblright{}.

\begin{cousinweb}
\textbf{\ur{شہد}} (\emph{shahd}) is a loan from Arabic \emph{shahd}, \textquotedblleft{}honey\textquotedblright{}, which reached Urdu through Persian. Its Arabic root is \emph{sh-h-d} (Lane).
\end{cousinweb}

\begin{grammarlens}[title={What you've built}]
One more word for this chapter.
\end{grammarlens}

\subsection*{Your turn}
\begin{itemize}
\raggedright
  \item \emph{Say it:} \emph{shahd}
  \item \emph{Say it:} \emph{shahd}, once more
  \item \emph{Say it:} say \emph{shahd}
  \item \emph{Recall:} say the Urdu for butter, then the Urdu for cheese, then say \emph{shahd} again
\end{itemize}

\begin{tcolorbox}[breakable,colback=teal!4,colframe=teal!35!black,title={Before you move on}]
What does \ur{شہد} mean? (\textquotedblleft{}Honey\textquotedblright{}.) Say it once more.
\end{tcolorbox}
\section[\texorpdfstring{\ur{انگور} (angūr)}{انگور (angūr)}]{\ur{انگور} (angūr) — grapes}

\label{lesson:UR-C56-angur}



\begin{quote}
Before the new one: say the Urdu for cheese, then the Urdu for honey.
\end{quote}

\subsection*{You'll want to know: \ur{انگور}}
\textbf{\ur{انگور}} — \emph{angūr} — \textquotedblleft{}grapes\textquotedblright{}.

\begin{cousinweb}
\textbf{\ur{انگور}} (\emph{angūr}) is a loan, usually traced to Persian \emph{angūr}, \textquotedblleft{}grape\textquotedblright{}.
\end{cousinweb}

\begin{grammarlens}[title={What you've built}]
One more word for this chapter.
\end{grammarlens}

\subsection*{Your turn}
\begin{itemize}
\raggedright
  \item \emph{Say it:} \emph{angūr}
  \item \emph{Say it:} \emph{angūr}, once more
  \item \emph{Say it:} say \emph{angūr}
  \item \emph{Recall:} say the Urdu for cheese, then the Urdu for honey, then say \emph{angūr} again
\end{itemize}

\begin{tcolorbox}[breakable,colback=teal!4,colframe=teal!35!black,title={Before you move on}]
What does \ur{انگور} mean? (\textquotedblleft{}Grapes\textquotedblright{}.) Say it once more.
\end{tcolorbox}
